\begin{table}[H]
\centering
\caption{Remuneración laboral por formalidad, 2010--2025}
\label{tab:remuneracion_geih_formalidad}
\footnotesize
\begin{tabular}{@{}p{1.6cm}p{7.0cm}rrr@{}}
\toprule
Grupo & Indicador & 2010 & 2025 & Crec. anual \\
\midrule
Formal & Remuneración laboral total mensual (Billones de pesos de 2025) & 13,7 & 29,3 & 5,20\% \\
Formal & Ocupados (Millones) & 5,4 & 10,1 & 4,33\% \\
Formal & Remuneración por trabajador (Millones de pesos de 2025 al mes) & 2,6 & 2,9 & 0,83\% \\
Formal & Horas semanales por trabajador & 50,0 & 46,1 & -0,54\% \\
Formal & Remuneración por hora trabajada (Miles de pesos de 2025) & 11,8 & 14,5 & 1,38\% \\
Informal & Remuneración laboral total mensual (Billones de pesos de 2025) & 11,2 & 13,4 & 1,22\% \\
Informal & Ocupados (Millones) & 11,2 & 12,5 & 0,71\% \\
Informal & Remuneración por trabajador (Millones de pesos de 2025 al mes) & 1,0 & 1,1 & 0,51\% \\
Informal & Horas semanales por trabajador & 45,5 & 42,0 & -0,53\% \\
Informal & Remuneración por hora trabajada (Miles de pesos de 2025) & 5,1 & 5,9 & 1,05\% \\
\bottomrule
\end{tabular}
\caption*{\footnotesize Nota: La remuneración por trabajador corresponde a una remuneración mensual equivalente en pesos constantes de 2025. La remuneración por hora se calcula como la remuneración laboral total dividida por las horas trabajadas. Fuente: cálculos propios con GEIH del DANE.}
\end{table}
